\subsection{営利組織の意思決定}

営利組織では、利益を生むことが重要な目標である。そのため、投資案の財務的な望ましさを評価する技法が使われる。

\subsubsection{許容最低収益率}

\begin{definitionbox}許容最低収益率\par\smallskip
許容最低収益率とは、組織が「投資として受け入れられる」と考える最低限の内部収益率である。英語では Minimum Acceptable Rate of Return、略して MARR である。
\end{definitionbox}

ある案の収益率が 10\% で、他の投資で確実に 20\% が期待できるなら、前者を選ぶ理由は弱い。MARR は機会費用を表す。ある案へ資金を投じることは、同じ資金を別の案へ投じる機会を捨てることでもある。

\subsubsection{経済寿命}

経済寿命とは、資産を使い続ける総費用が最小になる期間である。初期には、投資された資産の費用が大きい。一方、時間が経つほど運用費や保守費が増える。両者の合計が最小になる期間が、経済的に望ましい寿命である。

\subsubsection{計画対象期間}

提案の寿命が異なる場合、同じ期間で比較する必要がある。4 年の案と 6 年の案を比較するなら、どの期間を共通の分析期間とするかを決める。この共通期間が計画対象期間である。

\subsubsection{置換判断}

置換判断とは、既存資産を別の資産に置き換えるかどうかの判断である。レガシーソフトウェアを保守し続けるか、再開発するか、既製品へ移行するかといった判断が該当する。

置換判断では、埋没費用と残存価値を意識する。埋没費用は、すでに発生して回収できない費用である。経済的な判断では、過去に投じた費用への感情ではなく、将来の価値と費用を見なければならない。

\subsubsection{撤退判断}

撤退判断とは、ある活動から完全に離れる判断である。ソフトウェア製品の販売を終了する、特定サービスの提供をやめる、古いプラットフォームの対応を終了する、といった例がある。撤退は事前に計画されることもあれば、業績や技術状況の変化で急に必要になることもある。

\subsubsection{高度な営利判断の考慮事項}

判断の影響が大きい場合には、インフレーション、デフレーション、減価償却、税金なども考慮する必要がある。初学者は、金額だけでなく「会計上どう扱われるか」「時間とともに価値がどう変わるか」も意思決定に影響すると理解すればよい。

\par\smallskip
\begin{center}
\centering
\IfFileExists{assets/generated_figures/ch15/fig-ch15-04.png}{\includegraphics[width=0.96\linewidth]{assets/generated_figures/ch15/fig-ch15-04.png}}{\fbox{\parbox[c][48mm][c]{0.9\linewidth}{\centering 画像生成待ち\\営利組織の意思決定フロー}}}
\par\smallskip
\refstepcounter{figure}\label{fig:ch15-04}図 \thefigure: 営利組織の意思決定フロー
\end{center}
図 \thefigure は、営利組織の意思決定フローを視覚的に整理したものである。
\par\smallskip
